\section{Conclusion}\label{sec:conclusion}

Recorded sales are a censored measurement of demand, and models fit to raw sales inherit the censoring: they learn that stockout days sell nothing and capped days sell the cap, and then plan accordingly. The three case studies in this paper worked one repair each, selected by what the data records about the stockout. A binary on-shelf flag lets a smoothing recursion freeze its occurrence updates off shelf, so that a stockout run neither erodes the demand estimate nor drags the reopening forecast to zero. A per-observation cap lets a censored likelihood score capped days as lower bounds, which cut peak-day errors by half in our simulation precisely where the plain likelihood saturated at the cap. A fractional, noisily labeled availability supports a multiplicative saturating factor whose learned floor absorbs the label noise, fit here to $50{,}000$ series in minutes on a GPU. All three share one structure: the corruption is an explicit, switchable component of the model, and the demand forecast is the model run with the component switched off. And all three were evaluated under the same discipline: aggregate accuracy against recorded sales rewards reproducing the corruption, so demand models must instead be judged on calibration, stress-day behavior, and the plausibility of their counterfactuals. The mechanisms are simple, composable, and implemented in an open-source package; what remains genuinely open, endogenous availability, substitution, count likelihoods at scale, and the controlled cross-mechanism comparison, is catalogued in \S\ref{sec:discussion}.
